\begin{rubric}{Research Highlights}

\entry*[Scientific Software] Led the development of dozens of statistical and machine-learning software packages in R and Python (e.g. \textsc{SpectralUnmix}, \textsc{sagui}, \textsc{Capivara}, \textsc{Galmoss}, \textsc{qrpca}), deployed in galaxy surveys, IFU data analysis, and scalable astronomical data-processing pipelines.
%
\entry*[Statistical Inference] Co-developed the first likelihood-free inference framework applied to cosmology, establishing early simulator-driven inference workflows that anticipated modern simulation-based inference approaches in astrophysics.
%
\entry*[Astrostatistics] Co-authored invited review articles in the \emph{Annual Review of Statistics and Its Application} and the \emph{Notices of the American Mathematical Society}.
%
\entry*[Galactic Structure] Co-led the discovery of a previously unresolved high–pitch-angle structure in the Milky Way’s spiral arms, the identification of COIN--Gaia open clusters, and the construction of one of the largest catalogues of young stellar objects (SPICY) in the inner Galactic midplane.
%
\entry*[Nuclear Astrophysics] Developed the first hierarchical Bayesian framework for inferring nuclear reaction cross sections and S-factors in Big Bang nucleosynthesis, achieving state-of-the-art precision.
%
\entry*[Time-Domain Astronomy]
Contributed to the development and deployment of data-driven methods for time-domain astronomy, including alert-stream analysis, active-learning strategies, and broker-integrated science modules for ZTF and LSST science.

\end{rubric}
